\chapter{Einleitung}\label{chap:introduction}
    \section{Stand der Technik}\label{sec:state-of-the-art}
        \subsection{Additive Fertigungsverfahren im Überblick}\label{subsec:additive-manufacturing-overview}
            \Gls{additive-manufacturing}\index{Additive Fertigung} hat sich in den letzten Jahrzehnten vom reinen Rapid-Prototyping\index{Rapid Prototyping} zu Verfahren entwickelt, die auch in der Serienfertigung\index{Serienfertigung} eingesetzt werden \autocite[vgl.][S.~335--336]{urhalRobotAssistedAdditive2019}.
            Alle Verfahren zeichnen sich durch den schichtweisen Aufbau\index{Schichtaufbau} eines Bauteils auf Basis eines digitalen 3D-Modells\index{3D-Modell} aus.
            Dadurch lassen sich auch komplexe Geometrien fertigen, die mit klassischen, subtraktiven Verfahren\index{Subtraktive Fertigung} nur schwer oder nicht herstellbar wären \autocite[vgl.][S.~2785--2786]{nayyeriPlanarNonplanarSlicing2022}.

            Für die vorliegende Arbeit ist insbesondere die \gls{material-extrusion}\index{Materialextrusion} relevant, bei der ein thermoplastischer Kunststoff\index{Werkstoff!Thermoplast} aufgeschmolzen und über eine Düse\index{Düse} Schicht für Schicht abgelegt wird.
            Die Datenkette\index{Datenkette} aus \ac{CAD}-Modell\index{CAD-Modell}, \gls{slicing}\index{Slicing} und Maschinensteuerung\index{Maschinensteuerung} ist dabei in der Praxis weitgehend etabliert \autocite[vgl.][S.~3768--3770]{lettoriReviewGeometryRepresentation2022}.
            \sig{Adrian Gabriel Axtmann}

        \subsection{Konventionelle 3D-Drucksysteme}\label{subsec:conventional-3d-printing-systems}
            % TODO bei cartesian-kinematics n hinter Kartesische
            Die meisten heute verfügbaren 3D-Drucker\index{3D-Drucker!kartesisch} arbeiten mit einer \gls{cartesian-kinematics}\index{Kinematik!kartesisch}, d.\,h. mit drei zueinander senkrechten, translatorischen Achsen\index{Achse!translatorisch} \autocite[vgl.][S.~2787--2789]{nayyeriPlanarNonplanarSlicing2022}.
            Diese Bauform ist mechanisch vergleichsweise einfach, gut beherrschbar und wird durch eine breite Auswahl an Software unterstützt.
            Als Austauschformat\index{Datenformat!Austauschformat} für die Maschinensteuerung hat sich \ac{GCODE}\index{G-Code} durchgesetzt \autocite[vgl.][S.~2549--2551]{latifReviewCodeSTEP2021}.

            Die typische Prozesskette\index{Prozesskette} eines kartesischen Druckers ist linear aufgebaut: Anfangs wird das 3D-Modell aufbereitet und in horizontale Schichten\index{Schicht!horizontal} zerlegt.
            Daraus werden \glspl{toolpath} erzeugt und diese schließlich vom Drucker abgefahren \autocite[vgl.][S.~252--253]{gregoriSlicingTriangleMeshes2014}.
            Diese Vorgehensweise ist robust und einfach zu implementieren, weshalb sie sich als Quasi-Standard etabliert hat.
            Die meisten Slicer\index{Slicer} und Steuerungswerkzeuge sind jedoch stark auf diese drei Achsen ausgelegt, sodass eine Erweiterung auf \gls{nonplanar-slicing} oder auf mehrachsige Strategien\index{Schichtstrategie!mehrachsig} nur mit erheblichem Aufwand möglich ist \autocite[vgl.][S.~2--3]{pelzerAdditiveManufacturingNonplanar2021}.
            \sig{Adrian Gabriel Axtmann}

        \subsection{Limitationen von kartesischen 3D-Druckern}\label{subsec:limitations-cartesian-3d-printers}
            Aus der schichtweisen, planaren Bauweise\index{Bauweise!planar} ergeben sich mehrere bekannte Nachteile.
            Da das Material immer in horizontalen Ebenen\index{Schicht!horizontal} abgelegt wird, entstehen \gls{anisotropic-component-properties}\index{Anisotropie}: Bauteile sind in Z-Richtung deutlich weniger belastbar als in der XY-Ebene \autocite[vgl.][S.~1--2]{pelzerAdditiveManufacturingNonplanar2021}.
            Darüber hinaus müssen für \glspl{overhang}\index{Überhänge} oder freistehende Geometrien\index{Geometrie!freistehend} häufig \glspl{support-structures}\index{Stützstrukturen} mitgedruckt werden, welche nach dem Druck wieder entfernt werden müssen und zusätzlichen Material- und Zeitaufwand verursachen \autocite[vgl.][S.~537--538]{brandstoetterRoboterbasierter3DDruckMit2023}.

            Hinzu kommt, dass gekrümmte Oberflächen durch die treppenartige Schichtung sichtbare Stufen aufweisen (Treppeneffekt\index{Treppeneffekt}), was die Oberflächenqualität\index{Oberflächenqualität} beeinträchtigt \autocite[vgl.][S.~1--3]{ottonelloQuickCurveRevisitingSlightly2024}.
            Da Werkstück und Düse\index{Düse!Orientierung} nur in begrenzter Weise zueinander orientiert werden können, sind die Möglichkeiten der \gls{path-planning}\index{Bahnplanung} grundsätzlich auf zweidimensionale Konturen und \gls{infill}\index{Infill}-Muster beschränkt.
            Das räumliche Potenzial, das sich durch zusätzliche \gls{degrees-of-freedom}\index{Freiheitsgrad} ergeben würde, bleibt damit ungenutzt.
            \sig{Adrian Gabriel Axtmann}

        \subsection{Mehrachsige Fertigungssysteme}\label{subsec:multi-axis-systems-research-industry}
            Eine Möglichkeit, diese Einschränkungen zu umgehen, sind mehrachsige Fertigungssysteme\index{Mehrachsige Fertigung}.
            Insbesondere 6-Achs-Industrieroboter\index{6-Achs-Roboter} bieten zusätzlich zu den drei translatorischen\index{Freiheitsgrad!translatorisch} auch drei rotatorische \gls{degrees-of-freedom}\index{Freiheitsgrad!rotatorisch} und können so die Druckdüse\index{Druckdüse} nahezu beliebig im Raum orientieren \autocite[vgl.][S.~536--538]{brandstoetterRoboterbasierter3DDruckMit2023}.
            Damit lassen sich \gls{nonplanar-slicing}-Strategien\index{Schichtstrategie!nicht-planar} umsetzen, somit Schichten, die nicht zwingend horizontal verlaufen.
            In der Forschung gibt es hierzu bereits diverse Ansätze, von axialsymmetrischen Bahnstrategien\index{Bahnstrategie!axialsymmetrisch} \autocite[vgl.][S.~1--4]{lopez-arrabalAxisymmetricNonplanarSlicing2024} über zylindrische Vier-Achs-Verfahren\index{Bahnstrategie!zylindrisch}\index{4-Achs-Verfahren} \autocite[vgl.][S.~494--498]{guerraCylindricalSlicingAlgorithm2025} bis zur Kombination mit Lichtbogen- und Drahtprozessen \autocite[vgl.][S.~66--69]{dingAutomatedRoboticArcweldingbased2016}.

            In der industriellen Praxis\index{Industrielle Anwendung} sind solche Systeme aber bisher die Ausnahme.
            Das liegt weniger an der Hardware -- 6-Achs-Roboter sind ausgereift und etabliert -- sondern vielmehr an der Software.
            Die Kopplung von \gls{slicing} und \gls{inverse-kinematics}\index{Inverse Kinematik} ist nicht trivial \autocite[vgl.][S.~36--40]{aristidouInverseKinematicsTechniques2018}.
            Kollisionsfreie Bahnen\index{Bahnplanung!kollisionsfrei} müssen in oft engen Arbeitsräumen\index{Arbeitsraum} geplant werden \autocite[vgl.][S.~1--3]{bhattOptimizingPartPlacement2021}.
            Hinzu kommt, dass einheitliche Datenformate\index{Datenformat} fehlen, die alle benötigten Informationen (Geometrie, \gls{toolpath}, \gls{tool-orientation}\index{Werkzeugorientierung}) durchgängig abbilden \autocite[vgl.][S.~2549--2554]{latifReviewCodeSTEP2021}.

            An dieser Stelle, der softwareseitigen Verknüpfung von Slicer und Roboter, setzt die vorliegende Arbeit an.
            \sig{Adrian Gabriel Axtmann}

    \section{Grundidee der Arbeit}\label{sec:core-idea}
        \subsection{Übertragung auf den 6-Achs-3D-Druck}\label{subsec:transfer-to-6-axis-3d-printing}
            Die Grundidee dieser Arbeit ist es, das Prinzip des klassischen Slicens auf einen 6-Achs-Roboter zu übertragen.
            Statt das Bauteil ausschließlich in horizontale Schichten\index{Schicht!horizontal} zu zerlegen, sollen auch Schichten\index{Schicht!beliebig orientiert} mit beliebiger Orientierung im Raum betrachtet werden.
            Die etablierte planare Vorgehensweise wird dabei nicht verworfen, sondern als Spezialfall in ein erweitertes Modell eingebettet, das zusätzlich \gls{nonplanar-slicing} erlaubt.
            \sig{Adrian Gabriel Axtmann}

        \subsection{Konzeption eines 6-Achs-Slicer-Ansatzes}\label{subsec:concept-6-axis-slicer-approach}
            Auf dieser Idee aufbauend wird ein prototypischer \gls{6-axis-slicer}\index{6-Achs-Slicer} entworfen, der die Schichtzerlegung\index{Schichtzerlegung}, die \gls{path-planning} und die \gls{tool-orientation} gemeinsam betrachtet.
            Der Slicer bildet damit das Bindeglied\index{Toolchain} zwischen dem digitalen 3D-Modell und dem ausführbaren Programm für den Roboter.

            Die Software wird in zwei Module aufgeteilt:
            \textbf{\ac{medusa}}\index{Modul!Medusa} ist für das Laden und Anzeigen des 3D-Modells, das eigentliche Slicen sowie den Export der berechneten Pfade zuständig.
            \textbf{\ac{kronos}}\index{Modul!Kronos} nimmt diese Pfade entgegen und steuert über \ac{ros}\index{ROS} und \ac{moveit}\index{MoveIt} den Roboter an.
            Diese klare Trennung\index{Trennung!Planung/Ausführung} zwischen Planungs- und Ausführungsseite soll die Entwicklung und das Testen der einzelnen Komponenten erleichtern und die spätere Erweiterbarkeit\index{Erweiterbarkeit} verbessern.

            Im Rahmen der Arbeit werden insgesamt vier Slicing-Strategien betrachtet, die in unterschiedlichem Maße umgesetzt werden\index{Slicing-Strategien!Übersicht}:
            \begin{itemize}
                \item Der \textbf{planare Algorithmus}\index{Slicing-Strategien!planar}\index{Algorithmus!planar} dient als Referenz und Funktionsnachweis. Er stellt sicher, dass die gesamte Toolchain von der Modellaufbereitung bis zur Roboterbewegung grundsätzlich funktioniert.
                \item Der \textbf{Base-Algorithmus}\index{Slicing-Strategien!Base}\index{Algorithmus!Base} ist ein erster Multiachsen-Ansatz und ist im Rahmen dieser Arbeit ebenfalls vollständig implementiert.
                \item Der \textbf{Sphären-Algorithmus}\index{Slicing-Strategien!Sphären}\index{Algorithmus!Sphären} wird ausschließlich konzeptionell ausgearbeitet. Eine Implementierung ist nicht Bestandteil der Arbeit.
                \item Der \textbf{Kristall-Algorithmus}\index{Slicing-Strategien!Kristall}\index{Algorithmus!Kristall} wurde vollständig implementiert und bildet damit den nichtplanaren Slicing-Ansatz der Arbeit.
            \end{itemize}
            Diese gestaffelten Reifegrade\index{Reifegrad} der vier Strategien sind bewusst gewählt: Sie ermöglichen es, sowohl funktionierende Lösungen zu zeigen als auch die Grenzen einer prototypischen Umsetzung\index{Prototyp} ehrlich zu dokumentieren.
            \sig{Adrian Gabriel Axtmann}

        \subsection{Erwartete Vorteile gegenüber planaren Verfahren}\label{subsec:expected-benefits-over-established-approaches}
            Aus den in \cref{subsec:limitations-cartesian-3d-printers} beschriebenen Limitationen lassen sich umgekehrt die Potenziale eines 6-Achs-Drucks ableiten.
            Sie werden hier ausdrücklich als \emph{Erwartung} formuliert, deren Eintritt im Rahmen der prototypischen Umsetzung zu prüfen ist.

            Durch die freie \gls{tool-orientation}\index{Werkzeugorientierung!frei} könnten \glspl{overhang} teilweise ohne \glspl{support-structures} gefertigt werden.
            Da die Schichten der Bauteilgeometrie folgen können, ließe sich der Treppeneffekt an gekrümmten Oberflächen reduzieren.
            Zudem könnten die mechanischen \gls{anisotropic-component-properties} gezielter beeinflusst werden, indem die Schichten entlang der zu erwartenden Lastpfade\index{Lastpfad} ausgerichtet werden.
            Die zugehörigen Quellen sind bereits in \cref{subsec:limitations-cartesian-3d-printers} aufgeführt.

            Wie weit diese Potenziale in einem prototypischen Aufbau tatsächlich erreicht werden können, ist offen.
            Der Schwerpunkt der Arbeit liegt daher nicht darauf, diese Vorteile messtechnisch nachzuweisen, sondern die Voraussetzungen und Hürden zu identifizieren, die einer praktischen Realisierung im Weg stehen.
            \sig{Adrian Gabriel Axtmann}

    \section{Problemstellung}\label{sec:problem-statement-and-research-gap}
        \subsection{Herausforderungen bei der Bahnplanung}\label{subsec:technical-challenges-path-planning}
            Bei einem 6-Achs-Roboter müssen \gls{toolpath}, \gls{tool-orientation} und Bewegungsablauf gleichzeitig geplant werden.
            Die zusätzlichen \gls{degrees-of-freedom} vereinfachen die Aufgabe nicht, sondern vergrößern den Lösungsraum\index{Lösungsraum} und erhöhen die Komplexität der Randbedingungen\index{Randbedingung} \autocite[vgl.][S.~3776--3787]{lettoriReviewGeometryRepresentation2022}.

            Hinzu kommen prozessspezifische Anforderungen: Der Materialfluss\index{Materialfluss} aus der Düse, die Bahngeschwindigkeit\index{Bahnplanung!Geschwindigkeit} und die Orientierung müssen aufeinander abgestimmt sein, damit das Material gleichmäßig abgelegt wird \autocite[vgl.][S.~3--7]{pelzerAdditiveManufacturingNonplanar2021}.
            \sig{Adrian Gabriel Axtmann}

        \subsection{Offene Fragen bei Datenmodell und Prozesskette}\label{subsec:open-questions-data-model-and-process-chain}
            Die etablierten Datenformate und Werkzeugketten sind auf den planaren Druck zugeschnitten.
            Insbesondere \ac{GCODE}\index{G-Code!Beschränkung} ist auf drei Achsen ausgelegt und bildet zusätzliche Orientierungsinformationen nur über Erweiterungen oder Workarounds ab \autocite[vgl.][S.~2554--2560]{latifReviewCodeSTEP2021}.
            Für mehrachsige Anwendungen gibt es bislang kein einheitliches Format\index{Datenformat!fehlend}, das Geometrie, \gls{toolpath}, \gls{tool-orientation} und maschinenspezifische Parameter durchgängig zusammenführt \autocite[vgl.][S.~3--6]{witteRealisierungFunktionalDeterminierten2024}.

            Auch die Schnittstelle\index{Schnittstelle!Slicer-Controller} zwischen dem Slicer und dem Robotercontroller\index{Robotercontroller} ist nicht standardisiert.
            Verschiedene Controller verwenden unterschiedliche Bewegungsbefehle\index{Bewegungsbefehl} und Interpolationsverfahren\index{Interpolation}, sodass die Übertragung der geplanten Bahnen in eine reproduzierbare Roboterbewegung jeweils individuell gelöst werden muss \autocite[vgl.][S.~336--342]{urhalRobotAssistedAdditive2019}.
            \sig{Adrian Gabriel Axtmann}

        \subsection{Relevanz des Themas}\label{subsec:scientific-relevance-of-research-project}
            Wie in \cref{subsec:multi-axis-systems-research-industry} dargestellt, ist die softwareseitige Lücke\index{Software!Engpass} bei mehrachsiger additiver Fertigung der eigentliche Engpass.
            Sie wird umso spürbarer, je weiter der Standardanwendungsfall des planaren Drucks verlassen wird.
            Anwendungsfelder reichen vom großformatigen Druck\index{Großformatdruck} mit mehreren Robotern\index{Multi-Roboter-Fertigung} \autocite[vgl.][S.~1--4]{shenResearchLargescaleAdditive2019} bis hin zur medizinischen Implantatfertigung\index{Anwendung!Medizintechnik} \autocite[vgl.][S.~951--960]{martin-rodriguezRoboticPlatformNavigated2025}.

            Diese Arbeit greift einen kleinen Ausschnitt aus diesem Spektrum auf und untersucht prototypisch, wie sich ein Slicer für einen 6-Achs-Roboter umsetzen lässt und welche Konsequenzen sich daraus für eine spätere, ausgereiftere Umsetzung ergeben.
            \sig{Adrian Gabriel Axtmann}

    \section{Zielsetzung und Forschungsfragen}\label{sec:objectives-and-research-questions}
        \subsection{Gesamtziel der Arbeit}\label{subsec:overall-objective}
            % TODO hier auch wieder ändern: additiveN
            Ziel dieser Arbeit ist die prototypische Entwicklung\index{Prototyp} eines Slicers für \gls{nonplanar-slicing} in der robotergestützten \gls{additive-manufacturing} mit einem 6-Achs-Roboter.
            Der Slicer soll vom 3D-Modell ausgehend \glspl{toolpath} berechnen, die anschließend von einem Roboter ausgeführt werden können.

            Es geht dabei ausdrücklich nicht darum, ein produktionsreifes System zu entwickeln.
            Vielmehr soll gezeigt werden, dass eine durchgängige Kette\index{Prozesskette!durchgängig} von der Modell\-aufbereitung über das \gls{slicing} bis zur Roboteransteuerung grundsätzlich umsetzbar ist.
            Dabei sollen die Stellen identifiziert werden, an denen diese Kette in der Praxis an Grenzen stößt.
            Die Module \ac{medusa} und \ac{kronos} dienen dabei als prototypische Implementierung, welche die theoretischen Überlegungen greifbar machen und es erlauben, das Verhalten des Gesamtsystems zu beobachten.
            \sig{Adrian Gabriel Axtmann}

        \subsection{Forschungsfragen}\label{subsec:specific-research-questions}
            Auf Grundlage des Standes der Technik sowie der Problemstellung leiten sich die folgenden Forschungsfragen\index{Forschungsfragen} ab, welche die Arbeit strukturieren:

            \begin{enumerate}
                \item Welche Anforderungen muss ein Slicer erfüllen, um nicht-planare Bahnstrategien\index{Bahnstrategie!nicht-planar} für einen 6-Achs-Roboter zu erzeugen und wie unterscheidet er sich von einem klassischen planaren Slicer\index{Slicer!planar}?
                \item Wie kann ein solcher Slicer prototypisch umgesetzt werden und welche technischen Schwierigkeiten treten dabei auf?
                \item Welche Probleme entstehen bei der Übertragung der geplanten Bahnen auf einen realen Roboter (im konkreten Fall einen \ac{ur3e}\index{Roboter!UR3e} mit \ac{ros} und \ac{moveit})? Insbesondere: Wie lassen sich die beobachteten Effekte wie Konfigurationswechsel\index{Konfigurationswechsel} und Annäherungen an \glspl{singularity}\index{Singularität} einordnen?
                \item Welche Schlussfolgerungen lassen sich aus den umgesetzten und den nicht umgesetzten Slicing-Strategien für eine zukünftige Weiterentwicklung ziehen?
            \end{enumerate}

            Diese Forschungsfragen werden in den folgenden Kapiteln aufgegriffen und in der Diskussion zusammengeführt.
            \sig{Adrian Gabriel Axtmann}

        \subsection{Abgrenzung}\label{subsec:scope-delimitation}
            Die Arbeit beschränkt sich auf materialextrusionsbasierte Verfahren\index{Materialextrusion!Verfahren} und einen \ac{ur3e}-Roboter, der über \ac{ros} und \ac{moveit} angesteuert wird.
            Die beobachteten Effekte sind damit zunächst an diese konkrete Konfiguration gebunden, lassen sich aber teilweise auch auf andere 6-Achs-Knickarmroboter\index{Roboter!Knickarm} und vergleichbare Bewegungsplanungs-Frameworks\index{Bewegungsplanung!Framework} übertragen.

            \textbf{Nicht Bestandteil} der Arbeit sind
            \begin{itemize}
                \item hardwareseitige Entwicklungen am Roboter, am Extruder\index{Extruder} oder an der Sensorik\index{Sensorik},
                \item eine wirtschaftliche Bewertung\index{Wirtschaftlichkeit} des Verfahrens,
                \item ein mit realen Druckergebnissen validierter Slicer,
                \item ein quantitativer Vergleich der Druckqualität\index{Druckqualität} zwischen den vier Strategien.
            \end{itemize}

            Wie in \cref{subsec:concept-6-axis-slicer-approach} erläutert, sind der planare, der Base- und der Kristall-Algorithmus vollständig implementiert.
            Der Sphären-Algorithmus liegt ausschließlich als Konzept vor und wird im Rahmen dieser Arbeit nicht implementiert.
            \sig{Adrian Gabriel Axtmann}

    \section{Methodisches Vorgehen}\label{sec:methodological-approach}
        \subsection{Vorgehensweise}\label{subsec:research-design-and-development-strategy}
            Das Vorgehen ist folgt dem typischen Ablauf einer prototypischen Softwareentwicklung\index{Softwareentwicklung!prototypisch}: Anforderungsanalyse, Konzeption, Implementierung und Evaluation in mehreren Iterationen\index{Vorgehensmodell!iterativ}.

            Ausgangspunkt sind die Anforderungen, die sich aus dem Stand der Technik sowie der Problemstellung ergeben.
            Auf dieser Grundlage wird ein Konzept für die Software-Architektur\index{Softwarearchitektur} entwickelt, welches die in \cref{subsec:concept-6-axis-slicer-approach} eingeführten Komponenten \ac{medusa} und \ac{kronos} sowie deren Schnittstellen umfasst.
            Anschließend wird das Konzept implementiert und schrittweise erprobt.
            Erkenntnisse aus der Erprobung fließen in spätere Iterationen zurück, sodass Konzept und Implementierung sich gegenseitig schärfen.
            \sig{Adrian Gabriel Axtmann}

        \subsection{Evaluation}\label{subsec:evaluation-concept-and-comparison-criteria}
            Die Evaluation erfolgt in mehreren Schritten.
            Zunächst wird die Funktionsfähigkeit des Slicers anhand einfacher Testgeometrien\index{Testgeometrie} geprüft, um sicherzustellen, dass die berechneten \glspl{toolpath} plausibel sind.
            Daraufhin wird untersucht, inwieweit sich diese \glspl{toolpath} vom realen Roboter ausführen lassen, und welche Probleme dabei auftreten.
            Nachfolgend werden die umgesetzten Strategien konzeptionell mit dem klassischen, planaren Ansatz verglichen.

            Als Bewertungskriterien\index{Bewertungskriterium} dienen unter anderem die geometrische Korrektheit\index{Korrektheit!geometrisch} der \glspl{toolpath}, die Stabilität der Roboterbewegung\index{Roboterbewegung!Stabilität}, die Ausführbarkeit ohne Konfigurationswechsel oder Annäherung an \glspl{singularity} sowie qualitative Aspekte\index{Evaluation!qualitativ} wie \gls{support-structures}-Bedarf und \gls{toolpath}-Verlauf.
            Eine quantitative Vermessung\index{Evaluation!quantitativ} gefertigter Bauteile findet dagegen nicht statt.
            Eine derartige Untersuchung überschreitet den Rahmen der vorliegenden Arbeit und bleibt daher zukünftigen Forschungsprojekten vorbehalten
            \sig{Adrian Gabriel Axtmann}

    \section{Aufbau der Arbeit}\label{sec:thesis-structure}
        Die Arbeit ist in mehrere inhaltlich aufeinander aufbauende Kapitel gegliedert.

        Auf diese \textbf{Einleitung} folgen die \textbf{theoretischen Grundlagen}, in welchen für die Arbeit relevante Begriffe und Konzepte aus den Bereichen additive Fertigung, Robotik und Kinematik erläutert werden.
        Die \textbf{Anforderungsanalyse} leitet aus der Problemstellung konkrete funktionale und nicht-funktionale Anforderungen an das Gesamtsystem ab.
        Die \textbf{Konzeption} beschreibt die Architektur des Systems, die beiden Module \ac{medusa} und \ac{kronos}, die zwischen ihnen verwendete Schnittstelle sowie die vier betrachteten Slicing-Strategien.
        Die \textbf{Implementierung} dokumentiert die konkrete Umsetzung der Module und der Algorithmen.
        Die \textbf{Evaluation} zeigt die Ergebnisse der Erprobung anhand definierter Testfälle und beschreibt die dabei aufgetretenen Probleme bei der Roboteransteuerung.
        Die \textbf{Diskussion} ordnet die Ergebnisse ein, reflektiert die Grenzen der Arbeit (insbesondere die nicht erfolgreich umgesetzten Strategien und die Schwierigkeiten bei der Roboteransteuerung) und diskutiert mögliche Ursachen.
        Den Abschluss bildet das Kapitel \textbf{Fazit und Ausblick}, in dem die Ergebnisse zusammengefasst und konkrete Ansatzpunkte für eine Weiterentwicklung benannt werden.
        \sig{Adrian Gabriel Axtmann}

    \section{Begleitender Quellcode}\label{sec:accompanying-source-code}
        Ergänzend zu dieser schriftlichen Ausarbeitung ist der vollständige Quellcode\index{Quellcode} der Module \ac{medusa} und \ac{kronos} in einem GitLab-Repository\index{Repository!GitLab} verfügbar.
        Auf diese Weise lassen sich die in der \textbf{Implementierung} beschriebenen Konzepte direkt im Code nachvollziehen und die in der \textbf{Evaluation} dokumentierten Ergebnisse reproduzieren\index{Reproduzierbarkeit}.
        Die genaue Adresse des Repositorys sind im Anhang in \cref{sec:gitlab-repository-link} dokumentiert.
        \sig{Adrian Gabriel Axtmann}